\subsection{\bilingual{Proof of the Polynomial Decay of the Right Tail for \texorpdfstring{$a\in(0,1]$}{a∈(0,1]}}{在 \texorpdfstring{$a\in(0,1]$}{a∈(0,1]} 的情形下右尾部多项式衰减的证明}}

\begin{proof}
    \quad

    \bilingual{
        We first show the polynomial decay of the right-hand side for the case of $a\in(0,1)$. In light of \Cref{lemma:4.8}, we construct a killed branching random walk $\mathcal V^{\mathrm{k}}$ (killed according to a given rule specified later). Denote $\mathcal M^{\mathrm{k}}_n:=\min_{|u|=n}\mathcal V^{\mathrm{k}}(u)$ for $n\ge 0$. So long as $\mathcal V^{\mathrm{k}}$ satisfies the conditions listed in \Cref{lemma:4.8}, we have
    }{
        先证明 \eqref{eq:4.15} 在 $a\in(0,1)$ 的情形下成立. 根据\Cref{lemma:4.8}, 构造一个经灭杀的分支随机游走 $\mathcal V^{\mathrm{k}}$ (灭杀规则见后文具体说明). 对 $n\ge 0$ 记 $\mathcal M^{\mathrm{k}}_n:=\min_{|u|=n}\mathcal V^{\mathrm{k}}(u)$. 只要 $\mathcal V^{\mathrm{k}}$ 满足\Cref{lemma:4.8} 中列出的条件, 就有
    }
    \begin{equation*}
        \mathcal M^N_{\ell_N}
        \le
        \mathcal M^{\mathrm{k}}_{\ell_N}
        \quad
        \text{
            \bilingual{
                on the event $\{\tau>\ell_N\}$,
            }{
                在事件 $\{\tau>\ell_N\}$ 上,
            }
        }
    \end{equation*}
    \bilingual{
        where $\tau$ is the first time that the generation size of $\mathcal V^{\mathrm{k}}$ exceeds $N$. Note that the left-hand side in \eqref{eq:4.15}
    }{
        其中 $\tau$ 是 $\mathcal V^{\mathrm{k}}$ 的粒子总数首次超过 $N$ 的时刻. 注意到 \eqref{eq:4.15} 的左侧
    }
    \begin{align}
        &\quad
        \mathbb P\left(
            \mathcal M^N_{\ell_N}
            \ge
            \zeta'\log N
        \right)
        \nonumber
        \\
        &
        =
        \mathbb P\left(
            \mathcal M^N_{\ell_N}
            \ge
            \zeta'\log N,\,
            \tau\le\ell_N
        \right)
        +
        \mathbb P\left(
            \mathcal M^N_{\ell_N}
            \ge
            \zeta'\log N,\,
            \tau>\ell_N
        \right)
        \nonumber
        \\ &
        \le
        \mathbb P\left(
            \tau\le\ell_N
        \right)
        +
        \mathbb P\left(
            \mathcal M^{\mathrm{k}}_{\ell_N}
            \ge
            \zeta'\log N
        \right).
        \label{eq:prop:3.1:1}
    \end{align}
    \bilingual{
        Therefore, we need to construct a $\mathcal V^{\mathrm{k}}$ such that, with high probability,
    }{
        因此需要构造一个 $\mathcal V^{\mathrm{k}}$, 使得以高概率成立
    }
    \begin{itemize}
        \item
        \bilingual{
            $\tau>\ell_N$, i.e., the generation size of $\mathcal V^{\mathrm{k}}$ does not exceed $N$ till time $\ell_N$;
        }{
            $\tau>\ell_N$, 即 $\mathcal V^{\mathrm{k}}$ 的粒子总数在时刻 $\ell_N$ 之前不超过 $N$;
        }
        \item
        \bilingual{
            $\mathcal M^{\mathrm{k}}_{\ell_N}<\zeta'\log N$, i.e., the minimum position of $\mathcal V^{\mathrm{k}}$ is below $\zeta'\log N$ at time $\ell_N$;
        }{
            $\mathcal M^{\mathrm{k}}_{\ell_N}<\zeta'\log N$, 即 $\mathcal V^{\mathrm{k}}$ 的最左侧位置在时刻 $\ell_N$ 时低于 $\zeta'\log N$;
        }
    \end{itemize}
    \bilingual{
        and, the probability of deviating from these two normal situations, i.e., the right-hand side of \eqref{eq:prop:3.1:1}, decays at least polynomially with respect to $N$. Note that, the left-hand side of \eqref{eq:prop:3.1:1} is naturally hard to compute, since the selection mechanism underlying the $N$-BRW requires that we keep track of the whole population at each step, meaning that when we are determining whether a particular particle should be preserved, we have to look at the positions of all other particles, including the descendants of different ancestors at time $0$, which is hard to analyze. On the other hand, the right-hand side of \eqref{eq:prop:3.1:1} is easier to compute, since we can construct $\mathcal V^{\mathrm{k}}$ with a simpler selection mechanism, so that when we are determining whether a particular particle should be preserved, we only need to look at the position of itself, and thus we can analyze $\mathcal V^{\mathrm{k}}$ by using \nameref{lemma:many-to-one} and Mogulskii's estimate for random walks.
    }{
        并且, 偏离上述两种正常情况的概率, 即 \eqref{eq:prop:3.1:1} 的右侧, 至少以 $N$ 的多项式速率衰减. 注意到, \eqref{eq:prop:3.1:1} 的左侧自然难以计算, 因为 $N$-分支随机游走底层的选择机制要求我们在每一步都要跟踪整个种群, 意味着当我们在确定一个特定粒子是否应该被保留时, 我们必须查看所有其他粒子的位置, 包括时刻 $0$ 不同祖先的后代的位置, 这很难分析. 另一方面, \eqref{eq:prop:3.1:1} 的右侧更容易计算, 因为我们可以构造一个具有更简单选择机制的 $\mathcal V^{\mathrm{k}}$, 使得当我们在确定一个特定粒子是否应该被保留时, 我们只需要查看它自己的位置, 从而我们可以使用\nameref{lemma:many-to-one}和\nameref{theorem:mogulskii}对随机游走小偏差的估计来分析 $\mathcal V^{\mathrm{k}}$.
    }

    \bilingual{
        We first state our construction of the killed branching random walk $\mathcal V^{\mathrm{k}}$. Fix $c\in(0,t^{-1/3})$, $b\in(0,a)$, and $v>0$ such that
    }{
        首先给出经灭杀的分支随机游走 $\mathcal V^{\mathrm{k}}$ 的构造. 固定 $c\in(0,t^{-1/3})$, $b\in(0,a)$, 和 $v>0$ 使得
    }
    \begin{equation*}
        ct^{\frac 13}+(a-b)+v<1,\quad
        -b+\frac{\pi^2}2\sigma^2\left(c^{-2}t^{\frac 13}\right)+2v\left(1+t^{\frac 13}\right)<\zeta'.
    \end{equation*}
    \bilingual{
        Note that these inequalites are possible as $c\uparrow t^{-1/3}$, $b\uparrow a$, and $v\downarrow 0$. Also note that they are dependent on $a$, $t$, and $\zeta'$.
    }{
        注意到, 随着 $c\uparrow t^{-1/3}$, $b\uparrow a$, 和 $v\downarrow 0$, 这些不等式是可能满足的. 还要注意到它们依赖于 $a$, $t$, 和 $\zeta'$.
    }

\begin{figure}[H]
    \centering
    \begin{tikzpicture}[
        xscale=1.75,
        yscale=1.2,
        >=Stealth,
        % 样式定义
        survivor/.style={circle, fill=blue!70!black, inner sep=1.2pt},
        other-node/.style={circle, draw=blue!70!black, fill=white, inner sep=1pt},
        killed-node/.style={circle, draw=blue!70!black, fill=white, inner sep=1pt, cross out, thick, minimum size=3pt},
        edge-live/.style={thick, blue!70!black},
        edge-killed/.style={thick, dashed, blue!70!black},
        tube-fill/.style={blue!10, opacity=0.4},
        tube-edge/.style={thin, black},
        ref-dash/.style={thin, gray!60, dashed}
    ]

        % 1. 坐标轴
        \draw[->, thick] (0, -2.5) -- (0, 3.8) node[above] {$\mathcal{V}(u)$};
        \draw[->, thick, dashed] (0, 0) -- (6.5, 0) node[right] {$n$};

        % 2. 阶段 I: 水平管道 [0, 1.5]
        \fill[tube-fill] (0, -2.0) rectangle (1.5, 0);
        \draw[tube-edge] (0, -2.0) -- (1.5, -2.0) node[pos=0, left, black] {$-b \log N$};
        \draw[tube-edge] (0, 0) -- (1.5, 0);

        % 3. 阶段 II: 倾斜管道 [1.5, 5.0]
        \def\ystart{-1.8}
        \def\ylow{1.2}  % 终点下界
        \def\ywin{1.8}  % 终点窗口界
        \def\yhigh{2.8} % 终点上界
        \fill[tube-fill] (1.5, \ystart) -- (5.0, \ylow) -- (5.0, \yhigh) -- (1.5, \ystart+1.6) -- cycle;
        \draw[tube-edge] (1.5, \ystart) -- (5.0, \ylow);
        \draw[tube-edge] (1.5, \ystart+1.6) -- (5.0, \yhigh);

        \node[survivor] (y-node) at (1.5, \ystart) {};
        \node[left, xshift=-2pt] at (1.5, \ystart) {$y$};

        % 4. 均匀高频分支 (步长 0.3)
        \node[survivor] (n0) at (0, 0) {};

        % --- 阶段 I 路径 ---
        \foreach \i/\x/\y in {1/0.3/-0.4, 2/0.6/-0.1, 3/0.9/-0.8, 4/1.2/-1.4} {
            \coordinate (p\i) at (\x, \y);
            \node[survivor] at (p\i) {};
        }
        \draw[edge-live] (n0) -- (p1) -- (p2) -- (p3) -- (p4) -- (y-node);

        % 阶段 I 杀掉的分支 (管道外)
        \draw[edge-killed] (p1) -- (0.6, 0.4); \node[killed-node] at (0.6, 0.4) {};
        \draw[edge-killed] (p3) -- (1.2, -2.3); \node[killed-node] at (1.2, -2.3) {};
        \draw[edge-killed] (p4) -- (1.5, -1); \node[killed-node] at (1.5, -1) {};

        % --- 阶段 II 路径 (向上移动) ---
        \foreach \i/\x/\y in {1/1.8/-1.5, 2/2.1/-1.2, 3/2.4/-0.8, 4/2.7/-0.5, 5/3.0/-0.1, 6/3.3/0.3, 7/3.6/0.6, 8/3.9/0.9, 9/4.2/1.1, 10/4.5/1.3, 11/4.8/1.5} {
            \coordinate (q\i) at (\x, \y);
            \node[survivor] at (q\i) {};
        }
        \coordinate (m-node) at (5.0, 1.6);
        \draw[edge-live] (y-node) -- (q1) -- (q2) -- (q3) -- (q4) -- (q5) -- (q6) -- (q7) -- (q8) -- (q9) -- (q10) -- (q11) -- (m-node);
        \node[survivor] at (m-node) {};

        % 阶段 II 存活但非极小的分支 (管道内)
        \draw[edge-live] (q6) -- (3.6, 1.1) -- (3.9, 1.5) -- (4.2, 1.8) -- (4.5, 2.1) -- (4.8, 2.3);
        \foreach \x/\y in {3.6/1.1, 3.9/1.5, 4.2/1.8, 4.5/2.1, 4.8/2.3} \node[other-node] at (\x, \y) {};

        % 阶段 II 杀掉的分支 (越界)
        \draw[edge-killed] (q9) -- (4.5, 0.3); \node[killed-node] at (4.5, 0.3) {};
        \draw[edge-killed] (4.2, 1.8) -- (4.5, 2.8); \node[killed-node] at (4.5, 2.8) {};
        \draw[edge-killed] (4.8, 2.3) -- (5.0, 2.1); \node[killed-node] at (5.0, 2.1) {};

        % 5. 时间节点辅助线
        \draw[dashed, black!30] (1.5, -2.5) -- (1.5, 3.5) node[above, black] {$\ell$};
        \draw[dashed, black!30] (5.0, -2.5) -- (5.0, 3.5) node[above, black] {$\ell_N$};

        % 6. 标注窗口与跨度 (向右侧引出)

        % 左侧时刻 l 窗口
        \draw[ultra thick, blue!80!black] (1.5, -2.0) -- (1.5, -1.3);
        \draw[decorate, decoration={brace, amplitude=4pt, mirror}] (2, -2.0) -- (2, -1.3)
            node[midway, right, xshift=2pt] {$v \log N$};

        % --- 右侧引出线与标注 ---
        % 引出辅助水平线
        \draw[ref-dash] (5.0, \ylow) -- (6.2, \ylow);
        \draw[ref-dash] (5.0, \ywin) -- (5.8, \ywin);
        \draw[ref-dash] (5.0, \yhigh) -- (6.2, \yhigh);
        \draw[ref-dash] (1.53, -2.0) -- (2, -2.0);
        \draw[ref-dash] (1.53, -1.3) -- (2, -1.3);

        % 标注 M^k (紧贴 L_N 垂线右侧)
        \node[right, blue!70!black] at (5, 1.6) {$\mathcal{M}^{\mathrm{k}}_{\ell_N}$};

        % 标注 vL^1/3 (第一层大括号)
        \draw[ultra thick, blue!80!black] (5, \ylow) -- (5, \ywin);
        \draw[decorate, decoration={brace, amplitude=4pt, mirror}] (5.6, \ylow) -- (5.6, \ywin)
            node[midway, right, xshift=1pt] {$vL^{1/3}$};

        % 标注 cL^1/3 (第二层更远的大括号)
        \draw[decorate, decoration={brace, amplitude=5pt, mirror}] (6.3, \ylow) -- (6.3, \yhigh)
            node[midway, right, xshift=2pt] {$cL^{1/3}$};

        % 7. 阈值线
        \draw[red!80!black, thick, dashed] (0, 3.2) -- (5.5, 3.2) node[right] {$\zeta' \log N$};

    \end{tikzpicture}
    \caption{\bilingual{Illustration of the construction of the killed BRW coupled with the \texorpdfstring{$N$-BRW}{N-BRW} when \texorpdfstring{$a\in(0,1)$}{a in (0,1)}}{当 \texorpdfstring{$a\in(0,1)$}{a in (0,1)} 时与 \texorpdfstring{$N$-BRW}{N-BRW} 耦合的经灭杀的 BRW 的构造示意图}}
\end{figure}

    \begin{itemize}
        \item
        \bilingual{
            At time $0$, start with $\lfloor N^a\rfloor$ particles located at $0$; and they branch according to $\Xi$; note that these make $\mathcal V^{\mathrm{k}}$ satisfy the condition of \Cref{lemma:4.8};
        }{
            在时刻 $0$, 从 $\lfloor N^a\rfloor$ 个位于 $0$ 的粒子开始; 它们根据分支律 $\Xi$ 进行分支; 注意到这些设定使得 $\mathcal V^{\mathrm{k}}$ 满足\Cref{lemma:4.8} 的条件;
        }
        \item
        \bilingual{
            From time $0$ to $\ell:=\lfloor\log^{5/2}N\rfloor$, kill particles outside the tube $[-b\log N,0]$;
        }{
            从时刻 $0$ 到 $\ell:=\lfloor\log^{5/2}N\rfloor$, 灭杀管道 $[-b\log N,0]$ 以外的粒子;
        }
        \item
        \bilingual{
            At time $\ell$, further kill particles outside the interval $[-b\log N, (-b+v)\log N]$;
        }{
            在时刻 $\ell$, 进一步灭杀区间 $[-b\log N, (-b+v)\log N]$ 以外的粒子;
        }
        \item
        \bilingual{
            From time $\ell+1$ to $\ell_N$, for any alive particle in the last step at position $y$, kill particles outside a tilted tube with a small drift, i.e., $y+[r_j,s_j]$ for $j\in[\![\ell,\ell_N]\!]$, where $r_j:=\frac{\pi^2}{2}\sigma^2c^{-2}L^{-2/3}(j-\ell)$, $s_j:=r_j+cL^{1/3}$ and $L:=\ell_N-\ell$.
        }{
            从时刻 $\ell+1$ 到 $\ell_N$, 对于上一步中位置为 $y$ 的任何存活粒子, 灭杀一个带有小漂移的倾斜管道以外的粒子, 即 $y+[r_j,s_j]$ 对 $j\in[\![\ell,\ell_N]\!]$, 其中 $r_j:=\frac{\pi^2}{2}\sigma^2c^{-2}L^{-2/3}(j-\ell)$, $s_j:=r_j+cL^{1/3}$ 和 $L:=\ell_N-\ell$.
        }
        \item
        \bilingual{
            At time $\ell_N$, kill all particles outside the interval $y+[r_{\ell_N},r_{\ell_N}+vL^{1/3}]$ for any particle alive after the last step with an ancestor of time $\ell$ at position $y$.
        }{
            在时刻 $\ell_N$, 灭杀所有在上一步存活但在区间 $y+[r_{\ell_N},r_{\ell_N}+vL^{1/3}]$ 之外的粒子, 其中 $y$ 是其祖先在时刻 $\ell$ 的位置 .
        }
    \end{itemize}
    \bilingual{
        For formal analysis, we would denote the set of particles alive after each operation respectively as
    }{
        分别把每次操作后存活的粒子集合记为
    }
    \begin{equation*}
        \mathcal Z^{\mathrm{k}}_n,\,
        n\in[\![0,\ell]\!],\quad
        \mathcal Z^{\mathrm{k}}_{\ell+},\quad
        \mathcal Z^{\mathrm{k}}_n,\,
        n\in[\![\ell+1,\ell_N]\!],\quad
        \mathcal Z^{\mathrm{k}}_{\ell_N+}.
    \end{equation*}
    \bilingual{
        Note that
    }{
        注意到
    }
    \begin{equation*}
        \mathcal Z^{\mathrm{k}}_{\ell+}
        \subset
        \mathcal Z^{\mathrm{k}}_{\ell},\quad
        \mathcal Z^{\mathrm{k}}_{\ell_N+}
        \subset
        \mathcal Z^{\mathrm{k}}_{\ell_N}.
    \end{equation*}

    \bilingual{
        Next, we calculate an upper bound of the first term $\mathbb P(\tau\le\ell_N)$ in \eqref{eq:prop:3.1:1} for the constructed killed branching random walk $\mathcal V^{\mathrm{k}}$. Note that, by union bound and Markov's inequality,
    }{
        然后, 计算所构造的经灭杀的分支随机游走 $\mathcal V^{\mathrm{k}}$ 在 \eqref{eq:prop:3.1:1} 中的第一项 $\mathbb P(\tau\le\ell_N)$ 的上界. 注意到, 利用并集界和 Markov 不等式,
    }
    \begin{align*}
        \mathbb P(\tau\le\ell_N)
        &
        =
        \mathbb P\left(
            \begin{aligned}
                &\;\;\;\;\;
                (
                    \exists n\in[\![1,\ell-1]\!]:
                    \#\mathcal Z^{\mathrm{k}}_n>N
                )\\
                &\vee\;\;
                (
                    \#\mathcal Z^{\mathrm{k}}_{\ell+}>N
                )\\
                &\vee\;\;
                (
                    \exists n\in[\![\ell+1,\ell_N-1]\!]:
                    \#\mathcal Z^{\mathrm{k}}_n>N
                )\\
                &\vee\;\;
                (
                    \#\mathcal Z^{\mathrm{k}}_{\ell_N+}>N
                )
            \end{aligned}
        \right)
        \\ &
        \le
        \mathbb P\left(
            \exists n\in[\![1,\ell_N]\!]:
            \#\mathcal Z^{\mathrm{k}}_n>N
        \right)
        \\ &
        \le
        \sum_{n=1}^{\ell_N}
        \mathbb P\left(
            \#\mathcal Z^{\mathrm{k}}_n>N
        \right)
        % \\ &
        \le
        \sum_{n=1}^{\ell_N}
        \frac{
            \mathbb E\left[
                \#\mathcal Z^{\mathrm{k}}_n
            \right]
        }N.
    \end{align*}
    \bilingual{
        Let $\eta:=\max\{a,a-b+v+ct^{1/3}\}<1$. First we consider the operation from time $0$ to $\ell$, and the expected number $\mathbb E[\#\mathcal Z^{\mathrm{k}}_n]$ of particles alive till time $n\in[\![1,\ell]\!]$ is smaller than $N^{a+o_N(1)}$. Indeed, for any $n\in[\![1,\ell]\!]$, by \eqref{eq:4.10} in \Cref{lemma:4.7}, we have
    }{
        令 $\eta:=\max\{a,a-b+v+ct^{1/3}\}<1$. 首先考虑从时刻 $0$ 到 $\ell$ 的操作, 在时刻 $n\in[\![1,\ell]\!]$ 存活的粒子数量的期望 $\mathbb E[\#\mathcal Z^{\mathrm{k}}_n]$ 小于 $N^{a+o_N(1)}$. 事实上, 对于任意 $n\in[\![1,\ell]\!]$, 根据\Cref{lemma:4.7} 中的 \eqref{eq:4.10}, 有
    }
    \begin{align*}
        \mathbb E\left[
            \#\mathcal Z^{\mathrm{k}}_n
        \right]
        &
        =
        \mathbb E\left[
            \sum_{|u|=n}
            {\bf 1}_{
                \{
                    -b\log N
                    \le
                    \mathcal V^{\mathrm{k}}(u_j)
                    \le
                    0,\,
                    \forall j\in[\![1,n]\!]
                \}
            }
        \right]
        \\
        % &
        % =
        % \mathbb E\left[
        %     \sum_{i=1}^{\lfloor N^a\rfloor}
        %     \sum_{|u|=n}
        %     {\bf 1}_{
        %         \{
        %             -b\log N
        %             \le
        %             \mathcal V^{\mathrm{k}|(i)}(u_j)
        %             \le
        %             0,\,
        %             \forall j\in[\![1,n]\!]
        %         \}
        %     }
        % \right]
        % \\
        &
        =
        \sum_{i=1}^{\lfloor N^a\rfloor}
        \mathbb E\left[
            \sum_{|u|=n}
            {\bf 1}_{
                \{
                    -b\log N
                    \le
                    \mathcal V^{\mathrm{k}|(i)}(u_j)
                    \le
                    0,\,
                    \forall j\in[\![1,n]\!]
                \}
            }
        \right]
        \\
        &
        \le
        \sum_{i=1}^{\lfloor N^a\rfloor}
        \mathbb E\left[
            \sum_{n=1}^\infty
            \sum_{|u|=n}
            {\bf 1}_{
                \{
                    -b\log N
                    \le
                    \mathcal V^{\mathrm{k}|(i)}(u_j)
                    \le
                    0,\,
                    \forall j\in[\![1,n]\!]
                \}
            }
        \right]
        \\
        &
        =
        \lfloor N^a\rfloor
        O_N(\log N)
        =
        N^a
        O_N(\log N)
        =
        N^{a+o_N(1)}
        \le
        N^{\eta+o_N(1)},
    \end{align*}
    \bilingual{
        where, for the first equality, as defined, $\mathcal V^{\mathrm{k}}(u)=\mathcal V^{\mathrm{k}}(u,i)$ with $i$ the index of tree $\boldsymbol{\mathcal T}(i)$ which $u\in\mathbb V$ is on, and the summation is over all $i\in\mathbb I$ and $u\in\boldsymbol{\mathcal T}(i)$ with $|u|=n$. Then we consider the operation at time $\ell$, and the expected number $\mathbb E[\#\mathcal Z^{\mathrm{k}}_{\ell+}]$ of particles alive in the interval $[-b\log N, (-b+v)\log N]$ is of order $N^{a-b+v+o_N(1)}$. Indeed, by \eqref{eq:4.11} in \Cref{lemma:4.7}, we have
    }{
        其中, 第一个等式中右侧的求和项, 是对随机树族$\boldsymbol{\mathcal T}$上的点求和.
        % 其中,~ 对于第一个等式,~ 根据定义,~ $\mathcal V^{\mathrm{k}}(u)=\mathcal V^{\mathrm{k}}(u,i)$,~ 其中 $i$ 是 $u\in\mathbb V$ 所在的树 $\boldsymbol{\mathcal T}(i)$ 的索引,~ 并且求和是对所有 $i\in\mathbb I$ 和 $|u|=n$ 的 $u\in\boldsymbol{\mathcal T}(i)$ 进行的.
        然后, 考虑时刻 $\ell$, 在区间 $[-b\log N, (-b+v)\log N]$ 中存活的粒子数量的期望 $\mathbb E[\#\mathcal Z^{\mathrm{k}}_{\ell+}]$ 的数量级为 $N^{a-b+v+o_N(1)}$. 事实上, 根据\Cref{lemma:4.7} 中的 \eqref{eq:4.11}, 有
    }
    \begin{align*}
        \mathbb E\left[\#\mathcal Z^{\mathrm{k}}_{\ell+}\right]
        &
        =
        \mathbb E\left[
            \sum_{|u|=\ell}
            {\bf 1}_{
                \{
                    -b\log N
                    \le
                    \mathcal V^{\mathrm{k}}(u_j)
                    \le
                    0,\,
                    \forall j\in[\![1,\ell]\!]
                \}
            }
            {\bf 1}_{
                \{
                    \mathcal V^{\mathrm{k}}(u)
                    \le
                    (-b+v)\log N
                \}
            }
        \right]
        \\
        &
        =
        \mathbb E\left[
            \sum_{i=1}^{\lfloor N^a\rfloor}
            \sum_{|u|=\ell}
            {\bf 1}_{
                \{
                    -b\log N
                    \le
                    \mathcal V^{\mathrm{k}|(i)}(u_j)
                    \le
                    0,\,
                    \forall j\in[\![1,\ell]\!]
                \}
            }
            {\bf 1}_{
                \{
                    \mathcal V^{\mathrm{k}|(i)}(u)
                    \le
                    (-b+v)\log N
                \}
            }
        \right]
        \\
        &
        =
        \sum_{i=1}^{\lfloor N^a\rfloor}
        \mathbb E\left[
            \sum_{|u|=\ell}
            {\bf 1}_{
                \{
                    -b\log N
                    \le
                    \mathcal V^{\mathrm{k}|(i)}(u_j)
                    \le
                    0,\,
                    \forall j\in[\![1,\ell]\!]
                \}
            }
            {\bf 1}_{
                \{
                    \mathcal V^{\mathrm{k}|(i)}(u)
                    \le
                    (-b+v)\log N
                \}
            }
        \right]
        \\
        &
        =
        \lfloor N^a\rfloor
        N^{-b+v+o_N(1)}
        =
        N^{a-b+v+o_N(1)}.
    \end{align*}
    \bilingual{
        Now we consider the operation from time $\ell+1$ to $\ell_N$, and the expected number of particles alive till time $n\in[\![\ell+1,\ell_N]\!]$ is smaller than $N^{a-b+v+ct^{1/3}+o_N(1)}$. Indeed, by \eqref{eq:4.11} in \Cref{lemma:4.7}, \eqref{eq:4.8} in \Cref{lemma:4.6}, and the branching property, we have, for any $n\in[\![\ell+1,\ell_N]\!]$,
    }{
        现在考虑从时刻 $\ell+1$ 到 $\ell_N$ 的操作, 在时刻 $n\in[\![\ell+1,\ell_N]\!]$ 存活的粒子数量的期望 $\mathbb E[\#\mathcal Z^{\mathrm{k}}_n]$ 小于 $N^{a-b+v+ct^{1/3}+o_N(1)}$. 事实上, 根据\Cref{lemma:4.7} 中的 \eqref{eq:4.11}, \Cref{lemma:4.6} 中的 \eqref{eq:4.8} 和\nameref{lemma:branching-property}有, 对于任意 $n\in[\![\ell+1,\ell_N]\!]$,
    }
    \begin{align*}
        \mathbb E\left[\#\mathcal Z^{\mathrm{k}}_n\right]
        &
        =
        \mathbb E\left[
            \sum_{|u|=n}
            \begin{aligned}
                {\bf 1}_{
                    \{
                        -b\log N
                        \le
                        \mathcal V^{\mathrm{k}}(u_j)
                        \le
                        0,\,
                        \forall j\in[\![1,\ell]\!]
                    \}
                }
                {\bf 1}_{
                    \{
                        \mathcal V^{\mathrm{k}}(u_\ell)
                        \le
                        (-b+v)\log N
                    \}
                }
                \\[-5pt]
                {\bf 1}_{
                    \{
                        r_j
                        \le
                        \mathcal V^{\mathrm{k}}(u_j)
                        -
                        \mathcal V^{\mathrm{k}}(u_\ell)
                        \le s_j,\,
                        \forall j\in[\![\ell+1,n]\!]
                    \}
                }
            \end{aligned}
        \right]
        \\
        &
        =
        \sum_{i=1}^{\lfloor N^a\rfloor}
        \mathbb E\left[
            \sum_{|u|=n}
            \begin{aligned}
                {\bf 1}_{
                    \{
                        -b\log N
                        \le
                        \mathcal V^{\mathrm{k}|(i)}(u_j)
                        \le
                        0,\,
                        \forall j\in[\![1,\ell]\!]
                    \}
                }
                {\bf 1}_{
                    \{
                        \mathcal V^{\mathrm{k}|(i)}(u_\ell)
                        \le
                        (-b+v)\log N
                    \}
                }
                \\[-5pt]
                {\bf 1}_{
                    \{
                        r_j
                        \le\mathcal V^{\mathrm{k}|(i)}(u_j)-
                        \mathcal V^{\mathrm{k}|(i)}(u_\ell)
                        \le s_j,\,
                        \forall j\in[\![\ell+1,n]\!]
                    \}
                }
            \end{aligned}
        \right]
        \\
        &
        =
        \sum_{i=1}^{\lfloor N^a\rfloor}
        \mathbb E\left[
            \sum_{|\nu|=\ell}
            \sum_{
                |u|=n,\,
                u_\ell=\nu
            }
            \!\!\!\!\!\!\!\!\!
            \begin{aligned}
                {\bf 1}_{
                    \{
                        -b\log N
                        \le
                        \mathcal V^{\mathrm{k}|(i)}(\nu_j)
                        \le
                        0,\,
                        \forall j\in[\![1,\ell]\!]
                    \}
                }
                {\bf 1}_{
                    \{
                        \mathcal V^{\mathrm{k}|(i)}(\nu)
                        \le
                        (-b+v)\log N
                    \}
                }
                \\[-5pt]
                {\bf 1}_{
                    \{
                        r_j
                        \le\mathcal V^{\mathrm{k}|(i)}(u_j)-
                        \mathcal V^{\mathrm{k}|(i)}(u_\ell)
                        \le s_j,\,
                        \forall j\in[\![\ell+1,n]\!]
                    \}
                }
            \end{aligned}
        \right]
        \\ &
        =
        \sum_{i=1}^{\lfloor N^a\rfloor}
        \mathbb E\left[
            \begin{aligned}
                \sum_{|\nu|=\ell}
                {\bf 1}_{
                    \{
                        -b\log N
                        \le
                        \mathcal V^{\mathrm{k}|(i)}(\nu_j)
                        \le
                        0,\,
                        \forall j\in[\![1,\ell]\!]
                    \}
                }
                {\bf 1}_{
                    \{
                        \mathcal V^{\mathrm{k}|(i)}(\nu)
                        \le
                        (-b+v)\log N
                    \}
                }
                \\[-5pt]
                \times\;
                \mathbb E\left(
                    \sum_{
                        |u|=n,\,
                        u_\ell=\nu
                    }
                    \!\!\!\!\!
                    {\bf 1}_{
                        \{
                            r_j
                            \le\mathcal V^{\mathrm{k}|(i)}(u_j)-
                            \mathcal V^{\mathrm{k}|(i)}(u_\ell)
                            \le s_j,\,
                            \forall j\in[\![\ell+1,n]\!]
                        \}
                    }
                    \middle|
                    \mathcal F_\ell
                \right)
            \end{aligned}
        \right]
        \\ &
        =
        \sum_{i=1}^{\lfloor N^a\rfloor}
        \mathbb E\left[
            \begin{aligned}
                \sum_{|\nu|=\ell}
                {\bf 1}_{
                    \{
                        -b\log N
                        \le
                        \mathcal V^{\mathrm{k}|(i)}(\nu_j)
                        \le
                        0,\,
                        \forall j\in[\![1,\ell]\!]
                    \}
                }
                {\bf 1}_{
                    \{
                        \mathcal V^{\mathrm{k}|(i)}(\nu)
                        \le
                        (-b+v)\log N
                    \}
                }
                \\[-5pt]
                \times\;
                \mathbb E\left[
                    \sum_{|x|=n-\ell}
                    {\bf 1}_{
                        \{
                            r_{j+\ell}
                            \le\mathcal V(x_j)
                            \le s_{j+\ell},\,
                            \forall j\in[\![1,n-\ell]\!]
                        \}
                    }
                \right]
            \end{aligned}
        \right]
        \\ &
        \le
        \mathbb E\left[
            \sum_{|x|\le L}
            {\bf 1}_{
                \{
                    r_{j+\ell}
                    \le\mathcal V(x_j)
                    \le s_{j+\ell},\,
                    \forall j\in[\![1,L]\!]
                \}
            }
        \right]
        \\ &
        \quad\times\;
        \mathbb E\left[
            \sum_{|\nu|=\ell}
            {\bf 1}_{
                \{
                    -b\log N
                    \le
                    \mathcal V^{\mathrm{k}}(\nu_j)
                    \le
                    0,\,
                    \forall j\in[\![1,\ell]\!]
                \}
            }
            {\bf 1}_{
                \{
                    \mathcal V^{\mathrm{k}}(\nu)
                    \le
                    (-b+v)\log N
                \}
            }
        \right]
        \\ &
        =
        e^{(c+o_L(1))L^{1/3}}
        \mathbb E\left[
            \#\mathcal Z^{\mathrm{k}}_{\ell+}
        \right]
        =
        e^{(c+o_L(1))L^{1/3}}
        N^{a-b+v+o_N(1)}
        \\ &
        =
        e^{(c+o_N(1))(t^{1/3}\log N+o_N(\log N))}
        N^{a-b+v+o_N(1)}
        \\ &
        =
        N^{a-b+v+ct^{1/3}+o_N(1)}
        \le
        N^{\eta+o_N(1)},
    \end{align*}
    \bilingual{
        where we need to calculate the order of $L^{1/3}$ as
    }{
        其中需要计算 $L^{1/3}=(\ell_N-\lfloor\log^{5/2}N\rfloor)^{1/3}$ 的数量级为
    }
    \begin{align*}
        L^{1/3}
        % &
        % =
        % \left(
        %     \ell_N-\lfloor\log^{5/2}N\rfloor
        % \right)^{1/3}
        % \\ &
        =
        \left(
            (t+o_L(1))\log^3N+o_N(\log^3 N)
        \right)^{1/3}
        % \\ &
        =
        t^{1/3}\log N+o_N(\log N),
    \end{align*}
    \bilingual{
        and $\eta=\max\{a,a-b+v+ct^{1/3}\}<1$. We have shown that
    }{
        并且 $\eta=\max\{a,a-b+v+ct^{1/3}\}<1$. 现已证得
    }
    % \begin{equation*}
    $
        \mathbb E\left[
            \#\mathcal Z^{\mathrm{k}}_n
        \right]
        \le
        N^{\eta+o_N(1)},
        \,% \quad
        n\in[\![1,\ell_N]\!].
    $
    % \end{equation*}
    \bilingual{
        By union bound and Markov's inequality, we conclude that
    }{
        利用并集界和 Markov 不等式, 得出结论如下:
    }
    \begin{equation}
        \mathbb P\left(
            \tau\le \ell_N
        \right)
        \le
        \sum_{n=1}^{\ell_N}
        \frac{
            \mathbb E\left[
                \#\mathcal Z^{\mathrm{k}}_n
            \right]
        }{
            N
        }
        \le
        \ell_N
        N^{-(1-\eta)+o_N(1)}
        =
        N^{-(1-\eta)+o_N(1)}.
        \label{eq:4.19}
    \end{equation}
    \bilingual{
        We have calculated the small upper bound for the first term in \eqref{eq:prop:3.1:1}, the probability of the irregular event $\tau\le\ell_N$, and shown that this probability decays at least polynomially with respect to $N$.
    }{
        现已计算 \eqref{eq:prop:3.1:1} 中的第一项, 即 $\mathbb P(\tau\le\ell_N)$ 至少以多项式速率衰减.
    }

    \bilingual{
        Then we calculate a small upper bound of the second and last term in \eqref{eq:prop:3.1:1}, i.e., $\mathbb P(\mathcal M^{\mathrm{k}}_{\ell_N}\ge\zeta'\log N)$, for the constructed killed branching random walk $\mathcal V^{\mathrm{k}}$. Note that, for any $u\in\mathcal Z^{\mathrm{k}}_{\ell_N+}$,
    }{
        其次, 为所构造的经灭杀的分支随机游走 $\mathcal V^{\mathrm{k}}$, 计算 \eqref{eq:prop:3.1:1} 中的第二项也是最后一项, 即 $\mathbb P(\mathcal M^{\mathrm{k}}_{\ell_N}\ge\zeta'\log N)$ 的小上界, 注意到, 对于任意 $u\in\mathcal Z^{\mathrm{k}}_{\ell_N+}$,
    }
    \begin{align*}
        \mathcal V^{\mathrm{k}}(u)
        &
        \le
        \mathcal V^{\mathrm{k}}(u_\ell)
        +r_{\ell_N}
        +vL^{1/3}
        % \\ &
        \le
        (-b+v)\log N
        +\frac{\pi^2}2\sigma^2c^{-2}\ell_N^{1/3}
        +v\ell_N^{1/3}
        \\ &
        \le
        (-b+v)\log N
        +\frac{\pi^2}2\sigma^2c^{-2}(8t\log^3N)^{1/3}
        +v(8t\log^3N)^{1/3}
        \\ &
        \le
        \left(
            -b+\frac{\pi^2}2\sigma^2c^{-2}t^{1/3}+2v\left(1+t^{1/3}\right)
        \right)
        \log N
        % \\ &
        <\zeta'\log N,
    \end{align*}
    \bilingual{
        where $\ell_N\le 8t\log^3N$ holds for some $N_1=N_1(t)>0$, and all $N\ge N_1$. It means that any alive particle of the killed branching random walk $\mathcal V^{\mathrm{k}}$ at time $\ell_N$ is below the threshold $\zeta'\log N$. Hence,
    }{
        其中 $\ell_N\le 8t\log^3N$ 对于某个 $N_1=N_1(t)>0$ 和所有 $N\ge N_1$ 都成立. 这意味着 $\mathcal V^{\mathrm{k}}$ 在时刻 $\ell_N$ 的任何存活粒子都在阈值 $\zeta'\log N$ 以下. 因此,
    }
    \begin{align*}
        \mathbb P\left(
            \mathcal M^{\mathrm{k}}_{\ell_N}
            \ge
            \zeta'\log N
        \right)
        \le
        \mathbb P\left(
            \#\mathcal Z^{\mathrm{k}}_{\ell_N+}
            =
            0
        \right).
    \end{align*}
    \bilingual{
        We proceed to calculate a small upper bound for the probability of the irregular event $\mathcal Z^{\mathrm{k}}_{\ell_N+}=0$. Recall that $\mathbb E[\#\mathcal Z^{\mathrm{k}}_{\ell_+}]=N^{a-b+v+o_N(1)}$. Fix $\delta\in(0,\frac 12(a-b+v))$. Denote $\lfloor N^\delta\rfloor$ pairwise different elements in $\mathcal Z^{\mathrm{k}}_{\ell_+}$ as $\nu_i$, $i\in[\![1,\lfloor N^\delta\rfloor]\!]$. We have
    }{
        接着计算事件 $\mathcal Z^{\mathrm{k}}_{\ell_N+}=0$ 的概率的小上界. 回顾 $\mathbb E[\#\mathcal Z^{\mathrm{k}}_{\ell_+}]=N^{a-b+v+o_N(1)}$. 固定 $\delta\in(0,\frac 12(a-b+v))$. 把 $\mathcal Z^{\mathrm{k}}_{\ell_+}$ 中的 $\lfloor N^\delta\rfloor$ 个互异元素记为 $\nu_i$, $i\in[\![1,\lfloor N^\delta\rfloor]\!]$. 则有
    }
    \begin{align*}
        &\quad
        \mathbb P\left(
            \#\mathcal Z^{\mathrm{k}}_{\ell_N+}
            =
            0
        \right)
        \\
        &
        =
        \mathbb E\left[
            {\bf 1}_{
                \left\{
                    \#\mathcal Z^{\mathrm{k}}_{\ell_N+}
                    =
                    0
                \right\}
            }
            \left(
                {\bf 1}_{
                    \left\{
                        \#\mathcal Z^{\mathrm{k}}_{\ell_+}
                        \ge
                        N^\delta
                    \right\}
                }
                +
                {\bf 1}_{
                    \left\{
                        \#\mathcal Z^{\mathrm{k}}_{\ell_+}
                        <
                        N^\delta
                    \right\}
                }
            \right)
        \right]
        \\ &
        \le
        \mathbb E\left[
            {\bf 1}_{
                \left\{
                    \#\mathcal Z^{\mathrm{k}}_{\ell_N+}
                    =
                    0
                \right\}
            }
            {\bf 1}_{
                \left\{
                    \#\mathcal Z^{\mathrm{k}}_{\ell_+}
                    \ge
                    N^\delta
                \right\}
            }
        \right]
        +
        \mathbb P\left(
            \#\mathcal Z^{\mathrm{k}}_{\ell+}
            <
            N^\delta
        \right)
        \\ &
        =
        \mathbb E\left[
            \mathbb E\left(
                \prod_{
                    \nu\in
                    \mathcal Z^{\mathrm{k}}_{\ell+}
                }
                \!\!\!
                {\bf 1}_{
                    \left\{
                        \forall|u|=\ell_N,\,
                        u_\ell=\nu:\,
                        u\notin
                        \mathcal Z^{\mathrm{k}}_{\ell_N+}
                    \right\}
                }
                {\bf 1}_{
                    \left\{
                        \#\mathcal Z^{\mathrm{k}}_{\ell+}
                        \ge
                        N^\delta
                    \right\}
                }
                \middle|
                \mathscr F_\ell
            \right)
        \right]
        +
        \mathbb P\left(
            \#\mathcal Z^{\mathrm{k}}_{\ell+}
            <
            N^\delta
        \right)
        \\ &
        \le
        \mathbb E\left[
            \mathbb E\left(
                \prod_{i=1}^{
                    \lfloor N^\delta\rfloor
                }
                {\bf 1}_{
                    \left\{
                        \forall|u|=\ell_N,\,
                        u_\ell=\nu_i:\,
                        u\notin
                        \mathcal Z^{\mathrm{k}}_{\ell_N+}
                    \right\}
                }
                \middle|
                \mathscr F_\ell
            \right)
        \right]
        +
        \mathbb P\left(
            \#\mathcal Z^{\mathrm{k}}_{\ell+}
            <
            N^\delta
        \right)
        \\ &
        =
        \mathbb E\left[
            \prod_{i=1}^{
                \lfloor N^\delta\rfloor
            }
            \mathbb E\left(
                {\bf 1}_{
                    \left\{
                        \forall|u|=\ell_N,\,
                        u_\ell=\nu_i:\,
                        u\notin
                        \mathcal Z^{\mathrm{k}}_{\ell_N+}
                    \right\}
                }
                \middle|
                \mathscr F_\ell
            \right)
        \right]
        +
        \mathbb P\left(
            \#\mathcal Z^{\mathrm{k}}_{\ell+}
            <
            N^\delta
        \right),
    \end{align*}
    \bilingual{
        where we would proceed to use the branching property and \eqref{eq:4.7} in \Cref{lemma:4.6} to compute, for any $i\in[\![1,\lfloor N^\delta\rfloor]\!]$,
    }{
        其中继续利用分支性质和\Cref{lemma:4.6} 中的 \eqref{eq:4.7} 来计算, 对于任意 $i\in[\![1,\lfloor N^\delta\rfloor]\!]$,
    }
    \begin{align*}
        &\quad
        \mathbb E\left(
            {\bf 1}_{
                \left\{
                    \forall|u|=\ell_N,\,
                    u_\ell=\nu_i:\,
                    u\notin
                    \mathcal Z^{\mathrm{k}}_{\ell_N+}
                \right\}
            }
            \middle|
            \mathscr F_\ell
        \right)
        \\
        &
        =
        1-
        \mathbb E\left(
            {\bf 1}_{
                \left\{
                    \exists|u|=\ell_N:\,
                    u_\ell=\nu_i,\,
                    u\in
                    \mathcal Z^{\mathrm{k}}_{\ell_N+}
                \right\}
            }
            \middle|
            \mathscr F_\ell
        \right)
        \\ &
        =
        1-
        \mathbb E\left(
            {\bf 1}_{
                \left\{
                    \exists|u|=L:\,
                    u_\ell=\nu_i,\,
                    r_j
                    \le
                    \mathcal V^{\mathrm{k}}(u_j)
                    -
                    \mathcal V^{\mathrm{k}}(u_\ell)
                    \le
                    s_j,\,
                    \forall j\in[\![1,L]\!],\,
                    \mathcal V^{\mathrm{k}}(u)
                    -
                    \mathcal V^{\mathrm{k}}(u_\ell)
                    \le
                    r_{\ell_N}+vL^{1/3}
                \right\}
            }
            \middle|
            \mathscr F_\ell
        \right)
        \\ &
        =
        1-
        \mathbb E\left[
            {\bf 1}_{
                \left\{
                    \exists|x|=L:\,
                    r_{j+\ell}
                    \le
                    \mathcal V(x_j)
                    \le
                    s_{j+\ell},\,
                    \forall j\in[\![1,L]\!],\,
                    \mathcal V(x)
                    \le
                    r_{\ell_N}+vL^{1/3}
                \right\}
            }
        \right]
        \\ &
        =
        1-e^{-L^{1/3}o_L(1)}
        \le
        e^{
            -e^{
                -L^{\text{\fontsize{6pt}{5pt}$1/3$}}
                o_{\text{\fontsize{5pt}{5pt}$L$}}{\text{\fontsize{6pt}{5pt}$(1)$}}
            }
        }
        =
        e^{
            -N^{
                o_{\text{\fontsize{5pt}{5pt}$N$}}{\text{\fontsize{6pt}{5pt}$(1)$}}
            }
        },
    \end{align*}
    \bilingual{
        where the last inequality holds since $1-x\le e^{-x}$ for all $x\in[0,1]$, and the last equality holds since $L^{1/3}=t^{1/3}\log N+o_N(\log N)$ as computed above. Hence,
    }{
        其中最后一个不等式是因为对于所有 $x\in[0,1]$ 成立 $1-x\le e^{-x}$, 而最后一个等式是因为如上所计算的 $L^{1/3}=t^{1/3}\log N+o_N(\log N)$. 因此,
    }
    \begin{align*}
        \mathbb P\left(
            \#\mathcal Z^{\mathrm{k}}_{\ell_N+}
            =
            0
        \right)
        &
        \le
        \mathbb E\left[
            \prod_{i=1}^{
                \lfloor N^\delta\rfloor
            }
            e^{
                -N^{
                    o_{\text{\fontsize{5pt}{5pt}$N$}}{\text{\fontsize{6pt}{5pt}$(1)$}}
                }
            }
        \right]
        +
        \mathbb P\left(
            \#\mathcal Z^{\mathrm{k}}_{\ell+}
            <
            N^\delta
        \right)
        \\ &
        =
        e^{
            -N^{
                \delta+
                o_{\text{\fontsize{5pt}{5pt}$N$}}{\text{\fontsize{6pt}{5pt}$(1)$}}
            }
        }
        +
        \mathbb P\left(
            \#\mathcal Z^{\mathrm{k}}_{\ell+}
            <
            N^\delta
        \right).
    \end{align*}
    \bilingual{
        To calculate an upper bound for the second term in the last equation, we would use the fact that $\#\mathcal Z^{\mathrm{k}}_{\ell+}$ is stochastically larger than a binomial random variable $X$ and then apply Chernoff's bound for binomial distribution. Indeed, according to \eqref{eq:4.9} in \Cref{lemma:4.7}, we know that the probability for a branching random walk starting with one particle at $0$ to have at least one particle staying in the tube $[-b\log N,0]$ till time $\ell$ and ending up in the interval $[-b\log N, (-b+v)\log N]$ at time $\ell$ is of order $N^{-b+v+o_N(1)}$. Therefore, the number $\#\mathcal Z^{\mathrm{k}}_{\ell+}$ of particles alive after the same operations for $\mathcal V^{\mathrm{k}}$ starting with $\lfloor N^a\rfloor$ particles is (stochastically) larger than the number of particles at time $0$ to have at least one descendant alive after the same operations, which is a binomial random variable $X$ with parameters $\lfloor N^a\rfloor$ as the trial number and the aforementioned probability of order $N^{-b+v+o_N(1)}$ as the success probability, whose mean $\mathbb E[X]$ is of order $N^{a-b+v+o_N(1)}$. By Chernoff's bound for binomial distribution,
    }{
        为了计算上述等式中第二项的上界, 利用 $\#\mathcal Z^{\mathrm{k}}_{\ell+}$ 随机占优于一个二项分布随机变量 $X$ 的事实, 然后应用二项分布的 Chernoff 界. 事实上, 根据\Cref{lemma:4.7} 中的 \eqref{eq:4.9} 可知, 对于一个从$0$处以一个粒子开始的分支随机游走, 它至少有一个粒子在时刻 $\ell$ 前一直在管道 $[-b\log N,0]$ 内并且在时刻 $\ell$ 结束时在区间 $[-b\log N, (-b+v)\log N]$ 内的概率具有数量级 $N^{-b+v+o_N(1)}$. 因此, 对于从 $\lfloor N^a\rfloor$ 个粒子开始的$\mathcal V^{\mathrm{k}}$ , 在同样的操作下仍能存活的粒子数量 $\#\mathcal Z^{\mathrm{k}}_{\ell+}$, 随机占优于至少有一个后代在同样的操作下仍能存活的时刻 $0$ 处的粒子数量 (此因每一个这样的初始粒子最终可能贡献了不止一个存活的粒子), 后者也就是一个二项分布随机变量 $X$, 其试验次数为 $\lfloor N^a\rfloor$, 成功概率具有上述数量级 $N^{-b+v+o_N(1)}$, 其均值 $\mathbb E[X]$ 具有数量级 $N^{a-b+v+o_N(1)}$. 根据二项分布的 Chernoff 界,
    }
    \begin{align*}
        \mathbb P\left(
            \#\mathcal Z^{\mathrm{k}}_{\ell+}
            <
            N^\delta
        \right)
        &
        \le
        \mathbb P\left(
            \#\mathcal Z^{\mathrm{k}}_{\ell+}
            \le
            \lfloor N^\delta\rfloor
        \right)
        \\ &
        \le
        \mathbb P\left(
            X
            \le
            \lfloor N^\delta\rfloor
        \right)
        =
        \mathbb P\left(
            X
            \le
            \left(
                1
                -
                \lambda
            \right)
            \mathbb E[X]
        \right)
        \\ &
        \le
        e^{
            -\frac 12
            \lambda^2
            \mathbb E[X]
        }
        \le
        e^{
            -\frac 18
            \mathbb E[X]
        }
        \le
        e^{
            -N^{
                \delta+
                o_{\text{\fontsize{5pt}{5pt}$N$}}{\text{\fontsize{6pt}{5pt}$(1)$}}
            }
        },
    \end{align*}
    \bilingual{
        where
    }{
        其中
    }
    \begin{equation*}
        \lambda
        :=
        1
        -
        \frac{
            \lfloor N^\delta\rfloor
        }{
            \mathbb E[X]
        }
        \ge
        \frac 12,
        \quad
        N^\delta
        \le
        \frac 18
        N^{a-b+v}
    \end{equation*}
    \bilingual{
        holds for some $N_2=N_2(a,t,\zeta')>N_1$ and all $N\ge N_2$. We have shown that, for all $N\ge N_2$,
    }{
        对某个 $N_2=N_2(a,t,\zeta')>N_1$ 和所有 $N\ge N_2$ 成立. 现已证得, 对所有 $N\ge N_2$,
    }
    \begin{equation*}
        \mathbb P\left(
            \mathcal M^{\mathrm{k}}_{\ell_N}
            \ge
            \zeta'\log N
        \right)
        \le
        \mathbb P\left(
            \#\mathcal Z^{\mathrm{k}}_{\ell_N+}
            =
            0
        \right)
        \le
        2
        e^{
            -N^{
                \delta+
                o_{\text{\fontsize{5pt}{5pt}$N$}}{\text{\fontsize{6pt}{5pt}$(1)$}}
            }
        }
        =
        e^{
            -N^{
                \delta+
                o_{\text{\fontsize{5pt}{5pt}$N$}}{\text{\fontsize{6pt}{5pt}$(1)$}}
            }
        }.
    \end{equation*}
    \bilingual{
        Together with \eqref{eq:4.19} and \eqref{eq:prop:3.1:1}, we have, for some $\varepsilon=\varepsilon(a,t,\zeta')\in(0,\frac 12(1-\eta))$, some $N_0=N_0(a,t,\zeta')>N_2$, and all $N\ge N_0$,
    }{
        结合 \eqref{eq:4.19} 和 \eqref{eq:prop:3.1:1}可得, 对于某个 $\varepsilon=\varepsilon(a,t,\zeta')\in(0,\frac 12(1-\eta))$, 某个 $N_0=N_0(a,t,\zeta')>N_2$, 和所有 $N\ge N_0$,
    }
    \begin{align*}
        \mathbb P\left(
            \mathcal M^N_{\ell_N}
            \ge
            \zeta'\log N
        \right)
        &
        \le
        \mathbb P\left(
            \tau\le\ell_N
        \right)
        +
        \mathbb P\left(
            \mathcal M^{\mathrm{k}}_{\ell_N}
            \ge
            \zeta'\log N
        \right)
        \\ &
        \le
        N^{-(1-\eta)+o_N(1)}
        +
        e^{
            -N^{
                \delta+
                o_{\text{\fontsize{5pt}{5pt}$N$}}{\text{\fontsize{6pt}{5pt}$(1)$}}
            }
        }
        \le
        N^{-\varepsilon},
    \end{align*}
    \bilingual{
        which completes the proof of \eqref{eq:4.15} for the case $a\in(0,1)$.
    }{
        此即证得 $a\in(0,1)$ 情形下的 \eqref{eq:4.15}.
    }

    \bilingual{
        For $a=1$, fix any $b\in(0,1)$ sufficiently close to $1$ that $-b+\frac{\pi^2}{2}\sigma^2t<\zeta'$. Starting from fewer particles stochastically increases the minimum position. Applying the already proved $a=b$ case (with its corresponding $\varepsilon$ and $N_0$), for all $N\ge N_0$ we obtain
    }{
        对于 $a=1$, 取充分接近 $1$ 的 $b\in(0,1)$, 使 $-b+\frac{\pi^2}{2}\sigma^2t<\zeta'$. 从较少粒子出发会在随机序意义下增大最小位置. 将已证得的 $a=b$ 情形应用于此阈值 (其 $\varepsilon$ 和 $N_0$ 相应地选取), 对所有 $N\ge N_0$ 得到
    }
    \begin{equation*}
        \mathbb P\left(
            \mathcal M^{N|1}_{\ell_N}
            \ge
            \zeta'\log N
        \right)
        \le
        \mathbb P\left(
            \mathcal M^{N|b}_{\ell_N}
            \ge
            \zeta'\log N
        \right)
        \le
        N^{-\varepsilon}.
        \qedhere
    \end{equation*}
\end{proof}
